\documentclass[11pt]{article}

\usepackage[margin=1in]{geometry}
\usepackage{booktabs}
\usepackage{hyperref}
\usepackage{xcolor}
\usepackage{enumitem}
\usepackage{listings}
\usepackage{amsmath}
\usepackage{array}
\emergencystretch=3em

\hypersetup{
  colorlinks=true,
  linkcolor=blue,
  citecolor=blue,
  urlcolor=blue
}

\lstdefinestyle{code}{
  basicstyle=\ttfamily\small,
  breaklines=true,
  frame=single,
  columns=fullflexible,
  keepspaces=true
}
\lstset{style=code}

\title{Agentic Python Verifier:\\
Proof-Driven Python Code Generation via PlusCal Invariant Co-Synthesis}
\author{Harry Wang \and Eric Shang Nan Chen}
\date{Final Report\\May 28, 2026}

\begin{document}

\maketitle

\begin{abstract}
This project implements a prototype proof-driven Python generation pipeline.
Given a natural-language requirement, a synthesis agent proposes a PlusCal
algorithm, a candidate inductive invariant, a safety property, and finite
constant domains. A verifier translates the PlusCal module to TLA+ and checks
three obligations with TLC: initiation, consecution, and invariant implication
of the target property. When a check fails, the system parses TLC output into a
structured counterexample and sends that feedback to the synthesis agent for a
repair attempt. Only after all obligations pass does a refinement agent emit a
Python module with runtime assertions derived from the invariant. The final
repository contains the core pipeline, command-line interface, Anthropic
tool-call boundaries, TLC configuration generation, counterexample parsing,
example prompts, offline unit tests, mocked pipeline tests, and gated
integration tests for TLC and live end-to-end execution.
\end{abstract}

\section{Project Overview}

The goal of the Agentic Python Verifier is to invert the usual order of
large-language-model code generation. Instead of generating Python first and
then attempting to test or verify it afterward, the repository first builds a
formal artifact: a PlusCal algorithm with an explicit inductive invariant and
safety property. The verifier checks the formal artifact before the system is
allowed to refine it into executable Python.

The current implementation is organized around \texttt{src/main.py}. The
pipeline accepts a \texttt{TaskRequest}, constructs runtime settings, creates
artifact directories, asks the synthesis agent for an initial proposal, checks
the proposal with TLC, repairs failed proposals up to an iteration cap, and
then invokes the refinement agent only when the formal proof bundle passes.
The pipeline returns a structured \texttt{PipelineResult} with either
\texttt{verified} or \texttt{unverified} status.

The final repository includes:

\begin{itemize}[leftmargin=*]
  \item synthesis, repair, verification, and refinement agents;
  \item a formal layer for PlusCal translation, TLC invocation, TLC
  configuration generation, and TLC output parsing;
  \item Pydantic models for task requests, synthesis proposals, proof bundles,
  counterexamples, and pipeline results;
  \item a Typer command-line interface;
  \item example benchmark prompts for a bounded counter, bank transfer,
  resource semaphore, and ring buffer;
  \item offline tests that mock Anthropic and use captured TLC outputs; and
  \item gated integration tests for real TLC and live end-to-end execution.
\end{itemize}

\section{Motivation}

LLM-generated code is often plausible but difficult to trust. Unit tests can
catch many defects, but they sample behavior chosen by the developer or model.
For stateful algorithms, the most important requirements are often safety
properties such as ``the counter never exceeds the maximum,'' ``balances never
become negative,'' or ``the total amount of money is preserved.'' These
requirements naturally fit model checking, where all states in a finite model
can be explored rather than only selected examples.

The project therefore uses a formal intermediate representation as the trust
boundary. PlusCal gives an algorithmic notation that can be translated to
TLA+, and TLC can exhaustively check finite-state models \cite{lamport2002specifying,
lamport2009pluscal, yu1999model}. The language model still contributes
creativity: it proposes the algorithm, invariant, property, and repairs.
However, the verifier decides whether the proposal is acceptable. This creates
a separation of concerns: the model proposes, TLC checks, and Python is emitted
only after the formal artifact survives the proof obligations.

\section{Architecture}

The repository follows a staged architecture. The most important modules are
shown in Table~\ref{tab:architecture}.

\begin{table}[h]
\centering
\begin{tabular}{p{0.28\linewidth}p{0.62\linewidth}}
\toprule
Module & Responsibility \\
\midrule
\texttt{main.py} & Orchestrates the propose-check-repair-refine loop and
writes verified artifacts. \\
\texttt{synth\_agent.py} & Calls the LLM tool for initial PlusCal
and invariant proposals, then calls the repair tool after failures. \\
\texttt{verifier.py} & Runs PlusCal translation and all three TLC
proof obligations. \\
\texttt{refine\_agent.py} & Sends the verified module and invariant
conjuncts to the Python refinement tool. \\
\texttt{tlc\_config.py} & Generates TLC configuration files and
auxiliary TLA+ modules for consecution and property checking. \\
\texttt{counterexample\_parser.py} & Converts raw TLC output into
structured pass/fail/error/timeout results and counterexample traces. \\
\texttt{tools.py} & Defines the Anthropic tool schemas for proposal,
repair, and Python emission. \\
\bottomrule
\end{tabular}
\caption{Primary implementation components.}
\label{tab:architecture}
\end{table}

The complete dataflow is:

\begin{enumerate}[leftmargin=*]
  \item The user supplies a natural-language requirement.
  \item \texttt{SynthesisAgent.propose} requests a structured tool call
  containing a TLA+/PlusCal module, invariant name, property name, and finite
  constant domains.
  \item \texttt{Verifier.check} writes the proposal into a per-iteration work
  directory, invokes \texttt{pcal.trans}, writes auxiliary modules and
  \texttt{.cfg} files, and invokes TLC for each obligation.
  \item If any obligation fails, \texttt{SynthesisAgent.repair} serializes the
  failing proof bundle into an Anthropic \texttt{tool\_result} message.
  \item If all obligations pass, \texttt{RefineAgent.to\_python} requests a
  Python module and assertion map.
  \item The final TLA+ and Python files are written under the configured
  generated artifact directory.
\end{enumerate}

This architecture is intentionally narrow. The verifier has no LLM logic, the
LLM agents do not decide proof success, and the models make the boundary data
explicit.

\section{Formal Verification Approach}

For a generated algorithm, the verifier expects the translated TLA+ module to
define an initial predicate \texttt{Init}, transition relation \texttt{Next},
state tuple \texttt{vars}, candidate invariant \texttt{Inv}, and target
property \texttt{Property}. The verifier checks the three obligations in
Table~\ref{tab:obligations}.

\begin{table}[h]
\centering
\begin{tabular}{p{0.19\linewidth}p{0.25\linewidth}p{0.46\linewidth}}
\toprule
Obligation & Formula & TLC Encoding \\
\midrule
Initiation / reachable invariant check & $Init \Rightarrow Inv$ over the
generated \texttt{Spec} run &
Run TLC on the translated module with \texttt{SPECIFICATION Spec} and
\texttt{INVARIANT Inv}. This also checks reachable preservation under
\texttt{Spec}; a later reachable transition violation can appear in this run. \\
Consecution & $Inv \land Next \Rightarrow Inv'$ &
Generate \texttt{Consec\_<Module>.tla} with \texttt{IInit == Inv} and
an \texttt{ISpec} formula that starts from \texttt{IInit} and takes
\texttt{Next} steps over \texttt{vars}; check \texttt{Inv}. \\
Property implication & $Inv \Rightarrow Property$ &
Generate \texttt{Prop\_<Module>.tla} with \texttt{PInit == Inv} and an
unchanged transition; check \texttt{Property}. \\
\bottomrule
\end{tabular}
\caption{Implemented proof obligations.}
\label{tab:obligations}
\end{table}

The key design decision is to use \texttt{Inv} as an initial-state predicate
for the consecution and property checks. This makes TLC enumerate all states
that satisfy the candidate invariant and then verify either one-step
preservation or property satisfaction. The prompt therefore instructs the
synthesis model to make every variable enumerable with explicit finite set
membership, such as:

\begin{lstlisting}
Inv == /\ counter \in 0..MaxValue
       /\ pc \in {"Increment", "Done"}

Property == counter <= MaxValue
\end{lstlisting}

Without those finite membership constraints, TLC may be unable to enumerate
the invariant state space, or the search may be unbounded. The implementation
handles this by surfacing TLC evaluation errors and timeouts as structured
repair feedback rather than treating them as successful checks.

\section{Counterexample-Guided Repair Loop}

The repair loop is the main agentic component of the system. The synthesis
agent is not limited to one-shot generation. When a proposal fails, the
verifier returns a \texttt{ProofBundle} containing one result for initiation,
one for consecution, and one for property implication. Each
\texttt{ObligationResult} can represent \texttt{passed}, \texttt{failed},
\texttt{timeout}, or \texttt{error}.

For invariant and initial-state violations, the parser extracts the violated
predicate and a trace of TLC states. A consecution failure can therefore become
repair feedback in this shape:

\begin{lstlisting}
- consec: FAILED
    violated_predicate: Inv
    State 1 <Initial predicate>: counter=11
\end{lstlisting}

The repair agent receives this feedback together with the previous proposal
metadata. Its tool schema requires a revised module, the targeted obligation,
and a short diagnosis. This makes repair attempts inspectable: a later user or
test can see whether the model believed it was fixing an initiation,
consecution, or property implication failure.

The loop stops in one of two states. If all obligations pass before the
iteration cap, the system refines to Python. If the cap is exhausted, the
pipeline returns \texttt{unverified} and does not emit a Python module from the
failed proposal.

\section{Python Refinement}

Python generation is intentionally downstream of verification. The refinement
agent receives the verified TLA+/PlusCal module and a best-effort split of the
top-level invariant conjuncts. Its system prompt requires:

\begin{itemize}[leftmargin=*]
  \item a complete Python 3.11+ module using only the standard library;
  \item runtime assertions at the entry and exit of every state-mutating
  function;
  \item a module docstring containing the PlusCal module name, invariant, and
  property; and
  \item an \texttt{assertion\_map} connecting each TLA+ invariant conjunct to a
  Python check.
\end{itemize}

For example, an invariant over a bounded counter can be transformed into
Python checks around the mutating operation:

\begin{lstlisting}[language=Python]
def increment(counter: int, max_value: int = 10) -> int:
    assert 0 <= counter <= max_value
    if counter < max_value:
        counter += 1
    assert 0 <= counter <= max_value
    return counter
\end{lstlisting}

These assertions are not a substitute for proving the Python implementation
equivalent to the PlusCal algorithm. They are a refinement guard: the generated
code carries checks derived from the verified invariant, and the
\texttt{assertion\_map} provides traceability from the formal artifact to the
runtime code.

\section{Benchmarks and Examples}

The checked-in example prompts are intentionally small because TLC performs
finite-state exploration and because the main research question is whether the
pipeline structure works. The repository contains four example requirements:

\begin{lstlisting}
Implement a bounded counter from 0 to 10.
\end{lstlisting}

\begin{lstlisting}
Implement a bank transfer between two accounts that never allows a negative
balance and preserves the total sum of money.
\end{lstlisting}

\begin{lstlisting}
Implement a finite counting semaphore with Capacity = 2 and at most three
clients.
\end{lstlisting}

\begin{lstlisting}
Implement a finite ring buffer with Capacity = 3 over a small item set.
\end{lstlisting}

The test fixtures also include two bounded-counter TLA+ modules. One fixture
uses a suitable invariant for a counter bounded by \texttt{MaxValue}; the other
intentionally makes the invariant too strong by excluding the final allowed
counter value. The gated TLC integration test expects the good fixture to pass
all three obligations and the bad fixture to fail initiation or consecution
with a counterexample.

The current repository does not include saved successful live Anthropic
conversation logs or generated final Python artifacts from a real end-to-end
run. Therefore, this report treats the benchmarks as prompts and test fixtures,
not as evidence of completed live benchmark runs.

\section{Testing and Evaluation}

The project has a broad offline test suite. The tests avoid requiring an API
key or TLA+ tools jar by mocking the Anthropic boundary and using captured TLC
trace fixtures. The main tested areas are:

\begin{itemize}[leftmargin=*]
  \item model validation for task, synthesis, repair, and proof objects;
  \item runtime configuration and fail-fast error handling;
  \item CLI help text and missing-configuration behavior;
  \item TLC configuration and auxiliary module generation;
  \item TLC output parsing for passed runs, invariant violations, initial
  predicate violations, parse errors, and timeouts;
  \item Anthropic client overload fallback behavior;
  \item mocked synthesis and repair tool-call handling;
  \item mocked refinement and invariant-conjunct extraction; and
  \item mocked end-to-end pipeline behavior for first-try success, repair then
  success, and iteration exhaustion.
\end{itemize}

The repository also defines integration tests, but they are gated:

\begin{itemize}[leftmargin=*]
  \item \texttt{tests/test\_tlc\_runner.py} requires \texttt{TLA2TOOLS\_JAR}
  or \texttt{TLA\_TLC\_JAR}. It runs the verifier against known good and bad
  TLA+ fixtures.
  \item \texttt{tests/test\_e2e.py} requires both \texttt{ANTHROPIC\_API\_KEY}
  and a TLA+ tools jar. It performs a live bounded-counter pipeline run and
  may incur API costs.
\end{itemize}

This distinction is important for interpreting the final state. The checked-in
test suite demonstrates that the pipeline wiring, parser, configuration
generation, and mocked agent loop behave as intended. It does not, by itself,
prove that live Anthropic calls are deterministic or that the model reliably
solves every benchmark prompt. The code is designed so those live runs can be
performed, but this report does not claim live evaluation results that are not
present in the repository.

\section{Limitations}

The most important limitation is that verification is finite-domain model
checking, not unbounded theorem proving. The system asks the model to provide
small finite constant domains so TLC can explore the state space. This is
appropriate for a prototype but weaker than a general proof in TLAPS or another
theorem prover.

Second, the invariant must be enumerable as an initial-state predicate. This
is a practical constraint introduced by the consecution and property encodings.
An invariant that is logically reasonable but does not bound all variables
with finite membership constraints may cause TLC errors or timeouts.

Third, refinement is not formally verified. The PlusCal artifact can pass TLC,
and the generated Python can include invariant-derived assertions, but the
repository does not prove semantic equivalence between the PlusCal algorithm
and the Python implementation. The offline suite now includes deterministic
Python-level checks for a mocked generated module, but broader generated-code
testing remains future work.

Fourth, live LLM behavior is non-deterministic. The repository mitigates this
through structured tool schemas, prompt constraints, and offline mocked tests,
but repeatability ultimately depends on the checked TLA+ artifact, not on the
language-model transcript.

Finally, the invariant splitter in \texttt{RefineAgent} is a best-effort
top-level \texttt{/\\} scanner. It avoids splitting nested parenthesized,
bracketed, or braced expressions, but more complex TLA+ syntax may require a
syntax-aware parser to preserve structure reliably.

\section{Future Work}

The most immediate next step is to run and archive representative live
end-to-end evaluations with a real Anthropic API key and \texttt{tla2tools.jar}.
Those artifacts should include the prompt, final TLA+ module, TLC outputs,
repair iteration logs, and generated Python module.

Additional improvements include:

\begin{itemize}[leftmargin=*]
  \item expanding the benchmark suite beyond the current four prompt files to
  schedulers, locks, retry policies, and distributed protocol fragments;
  \item broadening Python-level tests for generated modules so more runtime
  assertions and functional behaviors are exercised after refinement;
  \item replacing the invariant-conjunct splitter with a TLA+-aware parser;
  \item enriching the current JSON/Markdown repair reports with links to TLC
  output excerpts and side-by-side proposal diffs;
  \item adding stronger validation of generated PlusCal before invoking TLC;
  and
  \item exploring a theorem-prover backend for unbounded proof obligations.
\end{itemize}

\section{Conclusion}

The final repository implements the core of a proof-driven Python generation
system. It can request a PlusCal/invariant/property proposal, translate it,
check initiation, consecution, and property implication with TLC, parse
counterexamples into structured repair feedback, and refine verified proposals
into Python with invariant-derived assertions. The offline test suite validates
the main internal contracts and the mocked agent loop, while gated integration
tests provide a path to real TLC and live end-to-end evaluation.

The project should be understood as a working prototype rather than a complete
verification platform. Its strongest contribution is the architecture: the LLM
is used for synthesis and repair, while TLC is responsible for accepting or
rejecting formal artifacts before Python code is produced.

\bibliographystyle{plain}
\bibliography{references}

\end{document}
